\section{Test de integraci\'on y speedup}
\label{sec:integracion}

El programa \texttt{tb\_integracion.asm} (anexo~\ref{ap:integracion}) ejercita en
una \'unica simulaci\'on todos los caminos relevantes de la FSM: relleno de las
dos v\'ias en los cuatro conjuntos, hits, escrituras hit que ensucian la v\'ia 0,
una escritura miss con write-around y cuatro reemplazos sucios consecutivos,
seguidos de accesos a MD Scratch y a registros IO.

\subsection{Cuenta esperada de eventos}

\begin{table}[htbp]
\centering
\begin{tabular}{lcccc}
\toprule
\textbf{Fase} & \sig{cont\_m} & \sig{cont\_w} & \sig{cont\_r} & \sig{cont\_cb} \\
\midrule
F0+F1: relleno v\'ia 0 (incluye constantes)         &  4 & 0 &  4 & 0 \\
F2: relleno v\'ia 1                                  &  4 & 0 &  4 & 0 \\
F3: 8 hits de lectura                                &  0 & 0 &  8 & 0 \\
F4: 4 hits de escritura sobre v\'ia 0                &  0 & 4 &  0 & 0 \\
F5: 1 write-around                                   &  1 & 0 &  0 & 0 \\
F6: 4 misses con reemplazo sucio                     &  4 & 0 &  4 & 4 \\
F7+F8: scratch e IO                                  &  0 & 0 &  0 & 0 \\
\midrule
\textbf{Total esperado}                              & \textbf{13} & \textbf{4} & \textbf{20} & \textbf{4} \\
\bottomrule
\end{tabular}
\caption{Reparto previsto de los contadores en \texttt{tb\_integracion.asm}.}
\label{tab:integracion-esperado}
\end{table}

Tras simular, anotar los valores reales: \dato{TBD}. Cualquier discrepancia respecto
a la tabla~\ref{tab:integracion-esperado} debe explicarse (orden FIFO, hits inesperados,
\ldots).

\subsection{Speedup frente a un sistema sin MC}

El speedup se define como
\begin{equation}
S = \frac{T_\textit{sin\_MC}}{T_\textit{con\_MC}}
\end{equation}
donde $T_\textit{con\_MC}$ es el n\'umero de ciclos medido en simulaci\'on hasta
entrar en el bucle final del programa, y $T_\textit{sin\_MC}$ es el coste estimado
si todo lw/sw cacheable fuese servido directamente por la MD palabra a palabra
(sin caching). Por construcci\'on, los accesos a scratch y a IO\_REG cuestan lo
mismo en ambos sistemas, por lo que se cancelan en la divisi\'on; basta con
restringir el c\'omputo a los lw/sw cacheable.

Sea $N_\textit{lw}$ y $N_\textit{sw}$ el n\'umero de lecturas y escrituras
cacheable, y $C_\textit{ovh}$ el coste de cada uno en el sistema sin caching:
$$
C_\textit{ovh} = T_\textit{arb} + 1 + \mathrm{CwW(MD)}
$$
(asumiendo que cada acceso es independiente; con burst no aplica porque cada lw
es una sola palabra). Entonces
\begin{equation}
T_\textit{sin\_MC} \approx \bigl(N_\textit{lw} + N_\textit{sw}\bigr)\cdot C_\textit{ovh}.
\end{equation}

Con $N_\textit{lw} = \mathrm{cont\_r}+\mathrm{cont\_m\_lw}=$ \dato{TBD},
$N_\textit{sw} = \mathrm{cont\_w}+N_\textit{wa} =$ \dato{TBD} y los costes de la
tabla~\ref{tab:latencias}:
$$
S = \frac{T_\textit{sin\_MC}}{T_\textit{con\_MC}} = \dato{TBD}.
$$

\paragraph{Lectura del cronograma.} Para obtener $T_\textit{con\_MC}$, basta marcar
en GTKWave el flanco de subida de \sig{reset} y el primer ciclo en que el PC del
MIPS queda atrapado en la instrucci\'on \texttt{beq R0, R0, -1} de F9 (el bucle
infinito). La diferencia en ciclos es $T_\textit{con\_MC}$.
